\subsection{Modelo autorregresivo logístico}\label{sec:autorregresivo}

Un modelo autorregresivo explica el valor presente de una serie a partir de
sus propios valores pasados. En su forma clásica, AR($p$), la observación
$y_t$ es una combinación lineal de los $p$ rezagos $y_{t-1},\dots,y_{t-p}$
más ruido. Aquí la respuesta es binaria (día de excedencia o no), así que la
dependencia temporal no puede entrar como media lineal; entra como término
del predictor de una regresión logística:
\begin{equation}
  \operatorname{logit} P(y_t = 1 \mid y_{t-1}, x_t)
    = x_t'\beta + \sum_{i=1}^{p} \alpha_i\, y_{t-i},
  \label{eq:arlog}
\end{equation}
donde $x_t$ reúne los tres puntajes factoriales del mismo día
(\cref{sec:analisis-factorial}) y los armónicos anual y semianual del
calendario. La \cref{sec:serie-tiempo} mostró que la autocorrelación de
rezago 1 sobrevive a la descomposición estacional; \cref{eq:arlog} es la
forma de aprovecharla en lugar de declararla como limitación, como se hizo en
\cref{sec:limitaciones-clasificacion}.

\subsubsection{Proceso}

El modelo se ajusta sobre la serie diaria de área metropolitana construida en
\cref{sec:serie-tiempo}, con $y_t = 1$ cuando al menos la mitad de las
estaciones que reportan excede el umbral de \texttt{excede\_anio\_5}. No se
modela la fracción de estaciones como binomial: la correlación transversal
produce una sobredispersión de 5.97, que subestimaría los errores estándar
por un factor de 2.4. El ajuste usa 2021--2024 y la validación 2025, la misma
partición que \cref{sec:clasificacion_tecs}. La especificación sale de cuatro
diagnósticos, no de una elección a mano:

\begin{enumerate}
  \item \textbf{Orden $p$.} La PACF de $y_t$ corta después del primer rezago
    (0.463 en el rezago 1, 0.100 en el 2) y el BIC, condicionado a las
    covariables, se minimiza en $p=1$. El rezago sobrevive a los factores
    contemporáneos ($LR = 71.4$, 1 g.l.): la persistencia no es un reflejo de
    que la meteorología sea persistente.
  \item \textbf{Estacionariedad.} Una serie binaria no puede tener raíz
    unitaria; lo que se verifica es homogeneidad temporal. El coeficiente
    autorregresivo no cambia entre 2021--2023 y 2024 ($p=0.64$), pero el
    nivel sí: una indicadora de 2024 en adelante es significativa
    ($p=0.0013$). Es un escalón que los factores no absorben y se declara
    como limitación.
  \item \textbf{Forma del predictor.} El \emph{linktest} rechaza sin
    curvatura y deja de rechazar ($p=0.60$) al incluir los cuadrados de
    ventilación y forzamiento fotoquímico, cuyos signos negativos indican
    saturación del efecto. El \emph{link} logit es correcto.
  \item \textbf{Dependencia residual.} Con el rezago, la autocorrelación de
    los residuales de Pearson en el rezago 1 cae de 0.212 a 0.025. Queda un
    residuo de 0.099 en el rezago 2 que no justifica un segundo rezago.
\end{enumerate}

\subsubsection{Resultados e interpretación}

Los coeficientes se leen como razones de momios, igual que en la
\cref{tab:momios-log}, y ningún intervalo del 95\,\% cruza el 1. Haber
excedido el día anterior multiplica los momios de exceder hoy por 3.83: es
el efecto que justifica el término autorregresivo. Entre los factores, el
forzamiento fotoquímico es el de mayor efecto (3.69 por desviación
estándar), seguido de la combustión primaria (2.37); la ventilación es el
único factor protector (0.32). Los cuadráticos quedan por debajo de 1 (0.70
y 0.69): el efecto marginal disminuye cuando el factor ya es alto.

\begin{figure}[t]
  \centering
  \includegraphics[width=\linewidth]{nowcast_roc.pdf}
  \caption{Curvas ROC en validación (2025) del modelo autorregresivo y de
    las dos líneas base. La persistencia produce tres vértices porque su
    predictor es binario.}
  \label{fig:nowcast-roc}
\end{figure}

Un AUC alto no basta; hay que compararlo contra lo que se obtendría sin
modelo. Las dos líneas base son la climatología (solo el calendario, AUC
0.668) y la persistencia (hoy será como ayer, AUC 0.717). El modelo alcanza
0.898 en 2025 y supera a ambas con $p<10^{-5}$ en la prueba de DeLong
(\cref{fig:nowcast-roc}). Con el umbral de Youden fijado sobre el ajuste
(0.327), no sobre la validación, la sensibilidad es de 0.759 y la
especificidad de 0.841: de 170 días de excedencia detecta 129, con 31 falsas
alarmas en 195 días limpios.

La validación cruzada dejando un año fuera separa el desempeño de la
partición del esperado en un año cualquiera: el AUC va de 0.840 (2021) a
0.915 (2024), con un agrupado de 0.878, IC 95\,\% $[0.859, 0.897]$. Esa
segunda cifra es la estimación de generalización; el 0.898 de 2025 cae
dentro del rango. El efecto de año no se explica por la prevalencia ni por la
cobertura de estaciones y es la limitación principal del modelo.

Dos precisiones acotan la lectura. El modelo subpredice en 2025 (media
predicha 0.338 contra 0.466 observado) por el escalón de nivel de 2024, y no
se recalibra porque sería ajustar un parámetro al periodo de prueba. Y es
\emph{nowcasting}, no pronóstico: los factores son del mismo día que la
respuesta, de modo que el modelo distingue días dentro de una temporada ---lo
que la climatología no puede hacer--- pero no anticipa el ozono de mañana.
